\subsection{The relative configuration sequence, with all maps specified}

Throughout this subsection, a finite $\Sigma$-group means a finite group
whose order has all prime divisors in $\Sigma$. Completion means the
inverse limit of all such quotients. A prime $r\in\Sigma$ is fixed.
For a complex smooth proper curve $C$ of genus $g\geq2$, put
\[
 \Lambda=\pi_1(C-P),\quad \Gamma=\pi_1((C\times C)-\operatorname{diag}),
 \quad B=\pi_1(C),\quad \Delta=B^{\Sigma}.
\]
The second projection is used, so that $\Lambda$ is its discrete normal
fibre subgroup. This exact discrete sequence is the
Fadell--Neuwirth sequence; the precise surface hypotheses and the
vanishing of $\pi_2(C)$ are given in
\citet[Section 2.1]{Bellingeri2004}. Superscripts $\Sigma$ below denote
pro-$\Sigma$ completion, not a quotient of the arithmetic Galois group.

\begin{claim}{A.CF1}{Proved}{The free-fibre monodromy lemma}
Let $1\to L\to H\to B_0\to1$ be an exact sequence of discrete groups,
where $L$ is free of finite rank at least two. If conjugation by $H$
acts through an $r$-group on $L/[L,L]L^r$, the natural continuous map
$L^{(r)}\to H^{(r)}$ is injective. Every prescribed finite $r$-quotient
of $L$ is consequently dominated by the restriction of a finite
$r$-quotient of $H$.
\end{claim}
\begin{proof}
Write $L_r=L^{(r)}$. Its iterated Frattini subgroups are open and
characteristic, and they are cofinal: their images eventually vanish in
each finite $r$-quotient. For a finite $r$-group $Q$, the kernel of
$\Aut(Q)\to\Aut(Q/\Phi(Q))$ is an $r$-group. Indeed, it acts freely
on the $|\Phi(Q)|^d$ ordered lifts of a fixed basis of $Q/\Phi(Q)$,
and every such lift generates $Q$. Passing to characteristic finite
quotients proves that the kernel of
\[
 \Aut(L_r)\longrightarrow\Aut(L_r/\Phi(L_r))
\]
is pro-$r$. Thus the closure of the conjugation image of $H$ in the
profinite group $\Aut(L_r)$ is pro-$r$. The action extends continuously
to $H^{(r)}$.

For completeness, $L_r$ has trivial centre. The Magnus map sends its
free generators to $1+X_i$ in the units of the noncommutative formal
power-series algebra $\F_r\langle\!\langle X_1,\ldots,X_d\rangle\!\rangle$.
It is injective: for each finite $r$-quotient $Q$ of $L_r$, substitution
$X_i\mapsto q_i-1$ is defined because the augmentation ideal of
$\F_r[Q]$ is nilpotent, and it recovers the quotient map to $Q$.
Nilpotence follows by induction on $|Q|$, using a central subgroup of
order $r$ and $(z-1)^r=0$. A power series commuting with both $X_1$
and $X_2$ is scalar. To see this, compare homogeneous terms: a
homogeneous polynomial commuting with $X_1$ has only the monomial
$X_1^n$, by comparing initial and terminal letters successively; then
commutation with $X_2$ eliminates every positive-degree term. A central
Magnus unit of augmentation one is therefore one.

The composite $L_r\to H^{(r)}\to\Aut(L_r)$ is the inner action,
so it is injective. A continuous injective map from a compact group
into a Hausdorff profinite group is a closed embedding. Given an open
normal subgroup $W$ of $L_r$, there is an open normal subgroup $N$ of
$H^{(r)}$ with $N\cap L_r\subset W$: use the induced topology and
then take the finite normal core of an open neighbourhood. The quotient
$H^{(r)}/N$ gives the stated domination.
\end{proof}

\begin{claim}{A.CF2}{Proved}{Exactness for a relative single-prime cusp kernel}
Set $F=\Lambda^\Sigma$, $E=\Gamma^\Sigma$, and
\[
 \begin{split}
 K_F&=\ker(F\to\Delta),\qquad
 K_E=\ker(E\to\Delta\times\Delta),\\
 F_{\crr}&=F/\ker(K_F\to K_F^{(r)}),\qquad
 E_{\crr}=E/\ker(K_E\to K_E^{(r)}).
 \end{split}
\]
The induced continuous sequence
\[
 1\longrightarrow F_{\crr}\longrightarrow E_{\crr}
       \xrightarrow{\pr_2}\Delta\longrightarrow1
\]
is exact. Every finite quotient of $F_{\crr}$ is dominated on the
fibre by a finite quotient of $E_{\crr}$ which has a normal $r$-subgroup
and a quotient $A\times A$, for a finite $\Sigma$-quotient $A$ of
$\Delta$.
\end{claim}
\begin{proof}
Let $q:F_{\crr}\twoheadrightarrow Q$ be continuous with $Q$ finite.
The image $R$ of $K_F^{(r)}$ is a normal $r$-subgroup, and
$A=Q/R$ is a quotient of $\Delta$. Let $a:C'\to C$ be the connected
Galois proper cover corresponding to this quotient. Put
\[
 Y=(C'\times C')-(a\times a)^{-1}(\operatorname{diag} C).
\]
Its removed divisor is the disjoint union of the graphs of the deck
transformations of $a$, so $Y\to C'$ has individually labelled
punctures. Choose a point of $C'$ above $P$. The fibre is
$V'=C'-a^{-1}(P)$. The surface fibration gives a discrete exact sequence
\[
 1\to L=\pi_1(V')\to H=\pi_1(Y)\to\pi_1(C')\to1.
\]
Here $L$ is free of rank $2g(C')+|A|-1\geq2$. In
$H_1(V',\F_r)$ the boundary classes span a subspace $W$, and there is
an exact sequence
\[
 0\to W\to H_1(V',\F_r)\to H_1(C',\F_r)\to0.
\]
Monodromy acts trivially on $W$, because the boundary sections are
labelled and oriented. It acts trivially on the quotient, because the
compactified family is the product $C'\times C'\to C'$. Choosing any
vector-space splitting shows that its image is contained in the additive
$r$-group $\Hom_{\F_r}(H_1(C',\F_r),W)$. This does not require
$r\nmid|A|$. By \statementref{A.CF1}, $L^{(r)}$ embeds in $H^{(r)}$.

The subgroup $H$ is normal in $\Gamma$, with quotient $A\times A$,
and $\Lambda\cap H=L$. The restriction of $q$ to $L$ takes values in
$R$. Choose a finite $r$-quotient of $H$, with kernel $J$, whose
restriction dominates this quotient of $L$. Thus
$J\cap L\subset\ker(q)$. Take
\[
 J^*=\bigcap_{\gamma\in\Gamma/H}\gamma J\gamma^{-1}.
\]
This is a finite intersection since $J$ is normal in $H$. The group
$H/J^*$ embeds in a finite product of $r$-groups. Consequently
$\Gamma/J^*$ is a finite $\Sigma$-group, an extension of $A\times A$
by an $r$-group, and $J^*\cap\Lambda\subset\ker(q)$. This last
containment remains true on completed fibre groups, as is seen in the
joint finite image of the two quotient maps.

The quotient of $\Gamma$ factors through $E_{\crr}$. Its image of
$K_E$ lies in $H/J^*$ because the $A\times A$ quotient comes from
the proper product; it therefore kills $\ker(K_E\to K_E^{(r)})$.
Conversely, the fibre map $F\to E_{\crr}$ has pro-$r$ image on
$K_F$, so it factors through $F_{\crr}$. Detection of every finite
quotient proves that this induced map is injective.

Finally, right exactness of pro-$\Sigma$ completion follows directly
from its universal property: quotienting $E$ by the closed image of
$F$ represents precisely the finite $\Sigma$-quotients of $B$, and
hence gives $\Delta$. The subgroup killed to form $E_{\crr}$ is
inside $K_E$, which is inside this fibre kernel. Thus the image of
$F_{\crr}$ is the whole kernel of $\pr_2$, and $\pr_2$ is onto.
\end{proof}

\begin{claim}{A.CF3}{Proved}{Specialization and the arithmetic fibre}
Let $C$ be smooth proper geometrically connected of genus at least two
over an algebraically closed field of characteristic $p$, and let
$\Sigma$ be any set of primes not containing $p$, with $r\in\Sigma$.
The preceding relative exact sequence holds for geometric etale groups.
If $C/k$ is defined over a finite field and $P\in C(k)$, its arithmetic
cuspidally pro-$r$ group is the inverse image of a chosen point section
$s_P(G_k)$ under the second projection of the arithmetic configuration
quotient. The comparison is over the identity of the whole $G_k$.
\end{claim}
\begin{proof}
We record the exact specialization input. SGA~1, Expose XIII,
Corollary~2.9, assumes a smooth proper morphism of finite presentation
$X\to S$, the complement $U$ of a relative normal-crossings divisor,
and the set of primes different from every residual characteristic of
$S$. For a locally constant constructible sheaf of finite groups with
these primes and tame along the boundary, it gives locally constant
$R^1$ and cohomological properness through degree one. Section~2.10
constructs the specialization homomorphism by strict localization and
a path; its ambiguity is inner conjugation. Applied to constant finite
groups, Corollary~2.9 identifies their torsors on the two geometric
fibres, so the resulting maximal prime-to-$p$ specialization map is an
isomorphism. It is functorial in the pair. Lemma~2.11 identifies the
images of inertia groups under this specialization. These statements
are in the opened text at pages 288--290, corresponding to original
margin pages 389--392 of \cite{SGA1}.

Here a suitable lift exists. SGA~1, Expose III, Theorem~7.3 assumes
$H^2(C,T_C)=H^2(C,\mathcal O_C)=0$ for a smooth projective special
fibre over the residue field of a complete local ring and gives a
smooth projective lift; Corollary~7.4 states the curve case. Both
vanishings hold for a curve. A chosen point lifts successively by
formal smoothness, and the compatible points give a section over the
complete ring. Equivalently, the marked deformation obstruction lies
in $H^2(C,T_C(-P))=0$. Take the Witt ring of a perfect algebraically
closed field of definition. The pairs
\[
 (\mathcal C,\mathcal P),\qquad
 (\mathcal C\times\mathcal C,\operatorname{diag}\mathcal C)
\]
are smooth proper pairs with smooth relative boundary. Their fibre
inclusion is transverse with multiplicity one. Functorial
specialization therefore identifies the groups and maps used in
\statementref{A.CF2}; the characteristic relative kernels and their
pro-$r$ quotients are identified as well.

For general algebraically closed fields, first descend the pointed
curve to a finitely generated field over $\F_p$, and take its
algebraic closure there. This is countable. Its Witt fraction field has
cardinality at most that of $\mathbb C$, so an algebraic closure
embeds into $\mathbb C$. Invariance of tame prime-to-$p$ groups of
smooth proper pairs under algebraically closed field extension, also
obtained from the same proper-pair comparison of finite torsors,
permits this descent and the return to the original field. Taking a
further maximal pro-$\Sigma$ quotient transfers the result for every
$\Sigma$ under consideration.

For the arithmetic assertion, write $\Pi_C$ for the geometrically
pro-$\Sigma$ group over the unchanged $G_k$. The configuration group
maps onto $\Pi_C\times_{G_k}\Pi_C$. Its kernel relative to this
product is geometric; replacing it by its pro-$r$ quotient gives a
group $\mathcal E_{\crr}$ with geometric subgroup $E_{\crr}$.
The actual fibre $C-P$ induces a map of extensions
\[
 \Pi_{C-P,\crr}\longrightarrow
 \mathcal E_{\crr}\times_{\Pi_C,s_P}G_k.
\]
On geometric kernels it is the isomorphism of
\statementref{A.CF2}, and on arithmetic quotients it is the identity
of $G_k$. Injectivity and surjectivity follow by the elementary
comparison of short exact sequences. Since $s_P$ is injective, the
right-hand group identifies with $\pr_2^{-1}(s_P(G_k))$.
Paths give exactly geometric inner conjugacy, not a change of the
arithmetic quotient.
\end{proof}

\begin{claim}{A.CF4}{Proved}{The common diagonal boundary and cusp transport}
In the arithmetic situation of \statementref{A.CF3}, the image
$\mathcal D$ of a fixed diagonal boundary component fits into a
continuous exact sequence
\[
 1\to\Z_r(1)\to\mathcal D\to\Pi_C\to1,
\]
and maps diagonally to $\Pi_C\times_{G_k}\Pi_C$. With paths chosen
through this boundary, the actual cusp subgroup of $C-P$ is
$\mathcal D\cap\pr_2^{-1}(s_P(G_k))$. Consequently, geometrically
conjugate proper point sections give an isomorphism of their actual
cuspidally pro-$r$ punctured groups over the identity of $\Pi_C$,
carrying the cusp subgroup to a geometric conjugate of the other cusp.
\end{claim}
\begin{proof}
Put $X=C\times C$, $D=\operatorname{diag} C$, and
$N=\mathcal O_X(D)|_D$. For an integer $n$ prime to $p$, the boundary
of $\sqrt[n]{(X,D)}$ is the gerbe of $n$-th roots of $N$ on $D$.
Here is the tame comparison being used, rather than a configuration
fibre exactness assertion. SGA~1 XIII, Proposition~5.2, assumes a
regular local scheme, a normal-crossings divisor with local equations
$f_i$, and a finite cover tame along that divisor. If $n_i$ is the
least common multiple of the relevant inertia orders, these integers
are invertible and the cover extends uniquely to a finite etale cover
after adjoining $T_i$ with $T_i^{n_i}=f_i$. The conclusion is over the
entire Kummer chart; the proof invokes purity. This is the opened
statement at pages 315--316, original margin page 429, of \cite{SGA1}.
Thus every finite tame cover of $X-D$ extends to a sufficiently high
root stack. Restriction to its boundary gerbe defines compatible maps
from the boundary fundamental groups to the configuration group.

A finite etale cover of the boundary gerbe with trivial stabilizer
action descends to $D$, and covers of $D$ pull back fully faithfully.
Etale locally the gerbe is $B\mu_n$; consequently the kernel of its
map to $\pi_1(D)$ is the image of $\mu_n$. An injection of this
finite stabilizer has not been assumed. Taking the inverse limit over
root orders, and then the geometrically pro-$\Sigma$ quotient, gives
a right exact sequence
\[
 \widehat\Z^{\Sigma}(1)\longrightarrow\mathcal T
                 \longrightarrow\Pi_C\longrightarrow1.
\]
Exactness follows also directly from the same trivial-stabilizer
descent criterion. The transition maps of the root gerbes are
surjective on fundamental groups: their maps on the base are the
identity and the maps of stabilizer images are onto. Compact inverse
limits preserve these right exact sequences. Right exactness of
pro-$\Sigma$ completion shows that discarding the geometric base
kernel introduces no further kernel over $\Pi_C$.

Let $\mathcal D$ be the image of $\mathcal T$ in
$\mathcal E_{\crr}$. Its map to the proper product is the injective
diagonal on $\Pi_C$, so its kernel over $\Pi_C$ is precisely the
image of inertia, a quotient of $\Z_r(1)$. It is the full
$\Z_r(1)$. Indeed, the geometric group of the fibre $C-P$ maps
onto the free pro-$r$ group on $2g$ generators. Its boundary loop is
the product of their commutators, up to orientation. Sending one pair
to $1+e_{12},1+e_{23}$ in the upper-unitriangular group over $\Z_r$
and the others to one sends that boundary loop to $1+e_{13}$, of
infinite $r$-adic order. The fibre injection of
\statementref{A.CF2} retains this inertia in the total quotient.
Transversality identifies it with diagonal inertia with multiplicity
one. This argument includes $r=2$.

The local transverse Kummer chart sends the actual cusp decomposition
group into $\mathcal D$ over $s_P(G_k)$, identically on inertia and
on the arithmetic residue group. Comparing the two short exact
sequences therefore identifies it with the whole inverse image of
$s_P(G_k)$ in $\mathcal D$. The paths for $P$ and $Q$ may both be
chosen in the same connected boundary tower, so this uses the same
$\mathcal D$ without assuming any prior separation of points.

Suppose $\gamma\in\Delta$ conjugates $s_P(G_k)$ to $s_Q(G_k)$.
Choose a geometric $d\in\mathcal D$ over $\gamma$. Conjugation by
$d$ identifies the two fibre subgroups and their cusp subgroups; on
the first proper coordinate it induces $\Inn(\gamma)$. The first
projection from the geometric part of the $Q$-fibre onto $\Delta$
is surjective. Choose there an $e$ mapping to $\gamma^{-1}$. Then
$\Inn(e)\Inn(d)$ is over the identity of $\Pi_C$ and $G_k$, and
carries the first cusp to a geometric conjugate of the second. No
lifting of a point section to a section of the configuration fibration
has been assumed.
\end{proof}

\subsection{Finite positive and negative tests in genus two}

\begin{claim}{A.CF5}{Proved}{A mixed-prime finite quotient detecting the relative fibre}
For genus two, $p=2$, $r=3$, and
$A=(\Z/3)^4\times\Z/5$, there is a quotient of the fibre of
order $1215$, with central cuspidal subgroup of order three and proper
quotient $A$, which embeds in a quotient of the configuration group of
order $492075$ with central kernel of order three over $A\times A$.
The defining surface braid relations and the fibre injection have been
checked by the retained program.
\end{claim}
\begin{proof}
Let $V=\F_3^4$ have standard alternating form
$J(v,w)=v_1w_2-v_2w_1+v_3w_4-v_4w_3$. On
\[
 G=(V\oplus V)\times\F_3\times\F_5^2
\]
define multiplication by adding the noncentral coordinates and adding
$2J(v_1+v_2,w_1+w_2)$ to the central coordinate. This is associative
because the added function is bilinear; $2$ is $1/2$ in $\F_3$.
Its order is $3^9\cdot25=492075$. Factor interchange preserves the
law. Form $G\rtimes\langle\tau\rangle$, where $\tau$ interchanges
the two $V$ coordinates and the two $\F_5$ coordinates.

Use \citet[Theorem 1.2]{Bellingeri2004} with $n=g=2$. Write its first
strand generators as $a,b,c,d$, put
$t=[a,b^{-1}][c,d^{-1}]$, and denote the interchange generator by
$\sigma$. Send $a,b,c,d$ to the four standard basis vectors in the
first $V$, giving $a$ also the coordinate $(1,0)$ in $\F_5^2$;
all central coordinates are zero. Send $\sigma$ to the product of
central coordinate $2$ with $\tau$. The eleven relations of the
presentation become identities: the four same-generator relations
use zero self-pairing, the four different-handle relations use zero
cross-handle pairing, the two same-handle mixed relations use the fact that
$[a,b]$ and $[c,d]$ both have central coordinate $1$ in $\F_3$,
and the surface relation
uses $-2=1$ in $\F_3$, giving $t=\sigma^2$.
The code checks each relation directly in this finite group.

The pure braid image lies in $G$ and is all of $G$. The first and
second strand images give eight independent $\F_3$ coordinates;
$a^3$ and its $\sigma$-conjugate give the two $\F_5$ coordinates,
and $[a,b]$ gives the central coordinate. The first fibre image is
\[
 Q=V\times\F_3\times\F_5\hookrightarrow G,
 \qquad (v,z,f)\longmapsto((v,0),z,(f,0)).
\]
This is injective by its coordinates and has order $3^5\cdot5=1215$.
The normal closure of $t$ is the central $\F_3$, in both groups.
Thus these finite maps factor through the respective cuspidally
pro-three groups over the full pro-odd proper groups, and their
proper quotients are exactly $A$ and $A\times A$. The program also
enumerates the $1215$ fibre elements and checks closure under its
specified generators. By \statementref{A.CF3}, the same geometric
finite test applies to the genus-two curve $y^2+y=x^5$ over
$\overline{\F}_2$.
\end{proof}

\begin{claim}{A.CF6}{Proved}{An explicit finite obstruction to full pro-odd exactness}
Let $F_4=\langle a,b,c,d\rangle$ be the first fibre in the genus-two
pure surface braid group $P_2$. Map $F_4$ onto the Heisenberg group
$H=\operatorname{UT}_3(\F_3)$ by sending $a,b$ to its standard
generators and $c,d$ to one. Let $K$ be the kernel and
$N=[K,K]K^5$. Then the finite odd quotient $F_4/N$, of order
$27\cdot5^{82}$, is not dominated on the fibre by any finite odd
quotient of $P_2$. In particular $F_4^{(2')}\to P_2^{(2')}$ is not
injective. The conclusion applies to the geometric groups of
$y^2+y=x^5$ in characteristic two by specialization.
\end{claim}
\begin{proof}
This is a finite certificate, supplemented by the proof that it
obstructs every possible finite odd total quotient. In Bellingeri's
presentation used above, set
$C=\sigma c\sigma^{-1}$ and $D=\sigma d\sigma^{-1}$.
Conjugation by them fixes $a,b$. On the other generators it is
\[
 \begin{array}{ll}
 \varphi_C(c)=tct^{-1},&
 \varphi_C(d)=tdt^{-1}(tc^{-1}tct^{-1}),\\
 \varphi_D(c)=ct^{-1},&
 \varphi_D(d)=tdt^{-1}.
 \end{array}                                                    \tag{CF.1}
\]
These are the actual conjugation automorphisms, not guessed
endomorphisms. For example the same-generator relation and its
$\sigma$-conjugate give
$CcC^{-1}=tct^{-1}$ and
$\varphi_C(t)=tc^{-1}tct^{-1}$. The mixed relation gives
$\varphi_C(dt^{-1})=tdt^{-1}$. Its conjugate gives
$\varphi_D(c)=ct^{-1}$; the analogous same-generator relation gives
the formula for $d$. The different-handle relations give the asserted
fixing of $a,b$. Normality of the discrete fibre follows from the
Fadell--Neuwirth sequence. The code independently verifies these
rewritten relations and the images of $t$ by free reduction.

The Heisenberg group is represented by triples with product
\[
 (x,y,z)(u,v,w)=(x+u,y+v,z+w+xv)\quad\text{over }\F_3.
\]
Both $\varphi_C^3$ and $\varphi_D^3$ preserve the quotient map to
$H$, as direct evaluation of all four generator images verifies.
Thus they normalize $K$, and they normalize $N$. A spanning tree in
the Cayley graph of this quotient gives a free Schreier basis of $K$
of rank $1+27(4-1)=82$. Its homology coordinates are computed by
recording oriented non-tree edges; the program verifies that its
basis-coordinate matrix is the identity.

Let $z=[a,b]\in H$. In the integral homology of $K$, write $c_s$
and $d_s$ for the lifted $c,d$ loops at vertex $s$. They are non-tree
edges, since $c,d$ map to one. On the independent vectors
\[
                 x=c_1-c_z,\qquad y=d_1-d_z
\]
the complete integral Schreier calculation gives
\[
 M_C=\begin{pmatrix}1&3\\0&1\end{pmatrix},\qquad
 M_D=\begin{pmatrix}1&0\\-3&1\end{pmatrix}.                    \tag{CF.2}
\]
Here $M_C$ is the action of $C^3$ and $M_D$ that of $D^3$.
The program exports the entire two $82$ by $82$ integral matrices
and checks all four equalities defining (CF.2), so this is not a
numerical approximation or a guessed invariant subspace. Hence
\[
 M_CM_D=\begin{pmatrix}-8&3\\-3&1\end{pmatrix},\qquad
 (M_CM_D)(x-y)=-(x-y)\pmod5.
\]
The class $x-y$ in $K/N$ is nonzero, having four nonzero Schreier
coordinates.

Suppose a finite odd quotient $P_2\to T$ with kernel $J$ dominated
$F_4/N$. Then $J\cap F_4\subset N$. Let
$H_0=\langle F_4,C^3,D^3\rangle$. It normalizes $K$ and $N$.
The quotient of the image of $K$ by the image of $N$ in $T$ is
exactly $K/N$, since $J\cap K\subset N$. The conjugation action of
$H_0$ on it therefore factors through the finite odd group
$\im(H_0\to T)$. In particular $C^3D^3$ must act with odd order.
Its action on the nonzero line generated by $x-y$ has order two,
a contradiction. No invariance of $K$ under the entire $P_2$ was
needed. Finally, an injective map between the two compact pro-odd
groups would induce the full topology on the fibre and would supply
a dominating finite total quotient for this specified finite fibre
quotient. Thus it cannot be injective.

The retained files are
\path{code/configuration_finite_tests.py}, its exact output, and
\path{code/configuration_finite_certificates.json}; the latter
contains the presentation inputs, the free words, the Schreier basis,
the integral matrices, and the sparse vector. The factorized arithmetic size is $27\cdot5^{82}$; its exact decimal
value is also retained in the program output.
\end{proof}

\subsection{Which sets of primes preserve the full configuration fibre?}

\begin{claim}{A.CF7}{Proved}{Exactness would pass to proper covers}
Suppose $C/\mathbb C$ has genus at least two and the full pro-$\Sigma$
fibre map for $(C\times C)-\operatorname{diag}$ is injective.
For every connected finite Galois etale $\Sigma$-cover $C'\to C$,
the analogous full fibre map for $C'$ is injective as well.
\end{claim}
\begin{proof}
First, if a subgroup $H$ of a discrete group $G$ contains a normal
subgroup of finite $\Sigma$-index, it has its full pro-$\Sigma$
topology induced from $G$. Given a normal $\Sigma$-index subgroup
of $H$, intersect with the fixed normal subgroup and take the
$G$-core. The quotient of the latter normal subgroup embeds in a
finite product of $\Sigma$-groups, and the resulting quotient of
$G$ is an extension of $\Sigma$-groups. It is therefore a
$\Sigma$-group. This proves the topology assertion and the
identification of the corresponding completed open subgroup.

Let $a:C'\to C$ have group $A$, and set
$Y=(C'\times C')-\bigcup_{\alpha\in A}\operatorname{graph}(\alpha)$.
It is a normal $A\times A$ cover of the original configuration
space. Choose $P'\in C'$ above $P$. Its fibre is
$V'=C'-a^{-1}(P)$, the corresponding proper-cover subgroup of
$C-P$. The preceding topology assertion and the assumed injection
identify $(\pi_1V')^\Sigma$ with a closed normal subgroup of
$(\pi_1Y)^\Sigma$; the quotient is $(\pi_1C')^\Sigma$.

Let $N$ be the closed normal subgroup of this fibre generated by
inertias at the points of $a^{-1}(P)$ other than $P'$. Monodromy
preserves these individual puncture labels. Thus the discrete total
group normalizes $N$, and so does its completion, since the
normalizer of a closed subgroup is closed. Filling the nonidentity
graphs in $Y$ kills exactly the closed total normal subgroup generated
by these same meridians. This elementary assertion follows by putting
a transverse disk through each smooth divisor and applying van
Kampen; every meridian along that divisor is a total conjugate of
its meridian in the chosen fibre. Since $N$ is already total-normal,
that total filling kernel equals $N$.

It follows that the quotient of the completed total group is
$(\pi_1((C'\times C')-\operatorname{diag}))^\Sigma$, and that its
injected fibre is $(\pi_1(C'-P'))^\Sigma$. Here we used right
exactness of pro-$\Sigma$ completion for both fillings; it follows
from the finite-quotient universal property exactly as in
\statementref{A.CF2}. This proves the assertion without assuming
that arbitrary completion commutes with kernels.
\end{proof}

\begin{claim}{A.CF8}{Proved}{A finite symplectic monodromy obstruction}
Let $C/\mathbb C$ have genus $h\geq2$, and let $r,\ell$ be distinct
primes. There is a connected degree-$r^3$ cover of $C-P$, with
$r$-group Galois group and cusp inertia of order $r$, whose compact
homology modulo $\ell$ has a subspace of dimension $2h-2$ on which
a subgroup of the configuration group preserving the cover acts
onto $\operatorname{Sp}_{2h-2}(\F_\ell)$. Consequently, if $r,\ell\in\Sigma$
and the full pro-$\Sigma$ configuration fibre were injective, every
prime divisor of $|\operatorname{Sp}_{2h-2}(\F_\ell)|$ would belong
to $\Sigma$.
\end{claim}
\begin{proof}
Choose a one-holed torus containing $P$ and a complementary connected
subsurface $M$ of genus $h-1$ with one boundary component. Map the
free group of $C-P$ onto $\operatorname{UT}_3(\F_r)$ by sending
the first two handle generators to its standard generators and all
generators supported on $M$ to one. Its boundary commutator is a
nontrivial central element $z$ of order $r$ (up to inversion).
Let $D^\circ\to C-P$ be the resulting regular cover and $D$ its
compactification. The restriction over $M$ is the disjoint union of
$r^3$ copies $M_s$, labelled by deck elements $s$. For a fixed $s$,
put
\[
 W=\{u_s-u_{zs}:u\in H_1(M,\F_\ell)\}\subset H_1(D,\F_\ell).
\]
The map defining $W$ is injective. Indeed, the usual symplectic handle
cycles in the disjoint subsurfaces $M_s$ have block-diagonal
intersection pairings with their respective dual cycles, so all their
classes are independent in $D$. This argument also works when
$\ell=2$, where the restricted pairing on $W$ itself can be
 degenerate. We use its identified vector space, not a false claim
that this restricted pairing is nondegenerate.

We give the topological calculation of the action explicitly. Let
$\alpha$ be an oriented simple nonseparating curve in $M$, joined to
$P$ by an embedded arc avoiding a fixed first-strand basepoint.
Pushing $P$ once round this based simple loop is the product of two
opposite Dehn twists on its left and right push-offs, denoted
$\alpha_0,\alpha_1$. To verify the identity, identify an annular
neighbourhood of the based loop with angular coordinate $\theta$ and
transverse coordinate $t$, and move the marked point by
$\theta\mapsto\theta+2\pi u$ during the isotopy parameter $u$.
At its two boundary collars the resulting map is one positive and
one negative twist. It is the identity outside this neighbourhood.
The two push-offs and the puncture bound a pair of pants. Their
cover monodromies are respectively $1,z$ (possibly after replacing
$z$ by $z^{-1}$); centrality makes the joining path immaterial.

The $r$-th power of each twist lifts and fixes the chosen sheet
outside the twisting annuli. On an annular cover of degree $n$,
where $n\mid r$, a twist of exponent $r$ lifts to exponent $r/n$:
in coordinates $\theta_{\rm base}=n\theta_{\rm cover}$ the lift is
$\theta_{\rm cover}\mapsto\theta_{\rm cover}+2\pi(r/n)f(t)$.
It is the identity on the boundary circles, so it glues to the
identity outside all annuli. This proves existence of the lift,
its sheet convention, and preservation of the cover without an
unproved assertion about outer actions.

The lifts of $\alpha_0$ have degree one and are the curves
$\alpha_s$ in the separate copies of $M$. Their twists have
exponent $r$. The lifts of $\alpha_1$ have degree $r$, one for
each central $\langle z\rangle$-orbit, and their twists have
exponent one. Each such long lift has homology class
\[
                        \sum_{j=0}^{r-1}[\alpha_{z^j s}]
                                                               \tag{CF.3}
\]
up to its common orientation. For this equality, lift the pair of
pants: each connected component is a degree-$r$ cyclic cover with
$r$ short boundaries, one long boundary, and one boundary above the
puncture. After compactification the last boundary bounds a disk.
The oriented boundary relation gives exactly (CF.3), with no extra
homology component.

The intersection of (CF.3) with $u_s-u_{zs}$ is zero: the two
coefficients are equal and cancel. Thus the long-lift twists act
trivially on $W$. The short-lift twists send
\[
 u_s-u_{zs}\longmapsto
 u_s-u_{zs}\ \mathbin{\pm}\ r\langle u,\alpha\rangle
                          (\alpha_s-\alpha_{zs}).              \tag{CF.4}
\]
The whole lifted point-push therefore preserves $W$ and induces the
usual symplectic transvection of exponent $\pm r$ on its identified
space $H_1(M,\F_\ell)$. All these lifts have the same sheet labels,
because they fix every point above the chosen basepoint outside the
annuli.

Since $r$ is invertible modulo $\ell$, powers give the transvection
of any coefficient. Every primitive integral homology class of $M$
is represented by a simple nonseparating curve: Euclidean handle
slides reduce a primitive vector to a handle basis vector, and
reversing the slides supplies the curve. Every nonzero vector modulo
$\ell$ has such a primitive integral lift (adjust one coordinate
by multiples of $\ell$ if necessary). We therefore obtain all
symplectic transvections on $H_1(M,\F_\ell)$.

These generate its full symplectic group. Here is a linear-algebra
proof, including characteristic two. A transvection of suitable
coefficient about $v-u$ sends $u$ to $v$ whenever
$\langle u,v\rangle\ne0$. If the pairing is zero, an intermediate
vector with both pairings nonzero exists, since two proper
hyperplanes do not exhaust the vector space. Thus the transvection
group is transitive on nonzero vectors. After fixing the first vector
$e$ of a symplectic basis, write its partner as $f+w+ae$, with $w$
in the symplectic complement. The product
$T_{e+w}T_w^{-1}T_e^{-1}$ fixes $e$ and removes $w$, with signs
chosen according to $\langle e,f\rangle=1$; a power of $T_e$
then removes $ae$. Induction in the symplectic complement fixes the
whole basis. This proves generation of $\operatorname{Sp}_{2h-2}$.

To obtain the last assertion, suppose the completed fibre is a
closed subgroup of the pro-$\Sigma$ configuration group. The
$r^3$-index subgroup defining $D^\circ$ has its full pro-$\Sigma$
topology there by the normal-core argument. Its quotient
$H_1(D,\F_\ell)$ is finite. The chosen point-push powers normalize
both the covering subgroup and the closed kernel defining this
homology quotient. Their closures still do so, because normalizers
of closed subgroups are closed. Conjugation on a finite quotient is
continuous, as one checks on finitely many coset representatives.
It therefore has finite $\Sigma$-image. Restricting to $W$ gives
a quotient of that image, but the preceding calculation makes this
quotient the full $\operatorname{Sp}_{2h-2}(\F_\ell)$. Its prime
divisors must lie in $\Sigma$.
\end{proof}

\begin{claim}{A.CF9}{Proved}{Classification of prime sets for the full fibre sequence}
For a smooth proper complex curve of genus at least two, the full
pro-$\Sigma$ configuration fibre sequence is exact if and only if
$\Sigma$ is empty, consists of one prime, or is the set of all
primes. If $C$ instead has characteristic $p$ and
$\Sigma\subseteq\{\ell:\ell\ne p\}$, it is exact if and only if
$|\Sigma|\leq1$. In particular the full prime-to-$p$ fibre map is
noninjective for every $p$ and every such curve, as claimed in R4.
The relative cuspidally pro-$r$ sequence of
\statementref{A.CF2} remains exact for every $r\in\Sigma$.
\end{claim}
\begin{proof}
The empty case is trivial. For $\Sigma=\{r\}$, monodromy on
$H_1(C-P,\F_r)=H_1(C,\F_r)$ is trivial, because the compactified
family is the product and there is only one boundary class, which
vanishes in abelian homology. Apply \statementref{A.CF1} and right
exactness. For the set of all primes, we use
\citet[Proposition 2.2]{LA_MT2008}, with a characteristic-zero
hyperbolic curve and the projection from its second to its first
configuration space: the proposition gives the full profinite fibre
exact sequence. Its prime-set hypothesis allows the total set (or a
singleton); it does not assert the sequence for a proper multi-prime
set.

Suppose now that $\Sigma$ is proper and contains distinct primes
$r,\ell$, and choose $p\notin\Sigma$. If the fibre sequence were
exact for $C$, it would be exact for every connected cyclic proper
cover of degree $r^m$, by \statementref{A.CF7}. Such covers exist by
the surjection $H_1(C,\Z)\twoheadrightarrow\Z/r^m$ and have genus
$h=1+r^m(g-1)$, hence arbitrarily large genus. Let $d$ be the
multiplicative order of $\ell$ modulo $p$ (for $p=2$, take $d=1$).
Choose $m$ such that $h-1\geq d$. Then $p$ divides
$\ell^{2d}-1$, and therefore divides
\[
 |\operatorname{Sp}_{2h-2}(\F_\ell)|
       =\ell^{(h-1)^2}\prod_{i=1}^{h-1}(\ell^{2i}-1).
\]
For reference, this order formula follows inductively by choosing a
nonzero first vector, its partner of pairing one, and a symplectic
basis of the complement; the successive factors are
$(\ell^{2i}-1)\ell^{2i-1}$. The divisibility contradicts
\statementref{A.CF8}, since $p\notin\Sigma$. Thus left injectivity
fails. Middle and right exactness still hold by the universal
finite-quotient argument, so this identifies the precise lost part
of the sequence.

For characteristic $p$, the functorial proper-pair specialization in
\statementref{A.CF3} identifies the pro-$\Sigma$ fibre map with the
complex one. As $p\notin\Sigma$, a nonempty such set is always
proper, giving exactly the stated classification. The mechanism also
explains the relative repair: a single retained cuspidal prime makes
all finite covering monodromy kernels used in
\statementref{A.CF2} $r$-groups; keeping two independent cusp primes
permits the symplectic finite actions above and their omitted-prime
divisors.
\end{proof}

\begin{claim}{A.CF10}{Proved}{The labelled several-section extension}
Let $C_0,B/\mathbb C$ be smooth proper curves, with $B$ of genus at
least one and $C_0$ of genus at least two, and let
$f_1,\ldots,f_m:B\to C_0$ have pairwise disjoint graphs. For
$Y=(C_0\times B)-\bigcup_i\operatorname{graph}(f_i)$, the relative
cuspidally pro-$r$ fibre sequence over pro-$\Sigma$ proper groups
is exact for every $r\in\Sigma$. The same assertion holds for simultaneously liftable
smooth proper pairs in characteristic prime to every member of
$\Sigma$ by the preceding specialization. In particular it applies
to the diagonal families after finite proper covers used in
cuspidally abelian reconstruction.
\end{claim}
\begin{proof}
For a finite fibre quotient, take its proper quotient $A$ and the
associated Galois cover $C_1\to C_0$. Pull this cover back by every
$f_i$, and take a common finite Galois $\Sigma$-cover $B_1\to B$
which trivializes all these $A$-torsors. Such a cover is given by
the intersection of their kernels; its group is a subgroup of a
finite product of copies of $A$. Now the inverse boundary in
$C_1\times B_1$ consists of disjoint labelled sections. Its compact
family is a product, its base is aspherical, and the homology action
modulo $r$ is the same shear group used in
\statementref{A.CF2}. Apply \statementref{A.CF1} and take the
normal core in $\pi_1(Y)$. The resulting total finite quotient is
an $r$-group by $A\times A_B$, where $A_B$ is the chosen base-cover
group; this is a $\Sigma$-group. Every finite fibre quotient is
detected, and the middle and right exactness proof is unchanged.
For the specific characteristic-$p$ diagonal-cover families, lift
the proper etale covers together with the smooth curve (finite
prime-to-$p$ etale covers lift uniquely on the proper-pair comparison)
and specialize. For a general characteristic-$p$ family of graphs,
the same conclusion requires a simultaneous lift of the indicated
smooth proper pair; the statement is understood with that hypothesis,
not as an assertion that arbitrary morphisms of curves lift.
\end{proof}

\begin{claim}{A.CF11}{Proved}{A finite two-prime test beyond the cofinite case}
For a genus-three curve, the full pro-$\{2,3\}$ configuration fibre
map is noninjective. A particular obstructing fibre quotient has order
$8\cdot3^{41}$; an exact computation of the preserving braid action
on a four-dimensional homology subspace gives
$\operatorname{Sp}_4(\F_3)$, of order $51840$, and an explicit
acting element of order five.
\end{claim}
\begin{proof}
Let $F_6$ have handle generators $a,b,c,d,e,f$, and let $K$ be the
kernel of its quotient onto $\operatorname{UT}_3(\F_2)$ obtained
from $a,b$ alone. Then $K$ is free of rank $1+8(6-1)=41$.
Use $N=[K,K]K^3$. The actual point-push lifts
$C,D,E,F$ are the $\sigma$-conjugates of $c,d,e,f$ in Bellingeri's
genus-three presentation. The program
\path{code/configuration_symplectic_crosscheck.py} derives their
conjugation words by the same relations as (CF.1), squares the words,
and checks preservation of the quotient to $\operatorname{UT}_3(\F_2)$.
It computes their action on the full integral Schreier homology and
on the submodule spanned by the four differences of the $c,d,e,f$
loops at the identity and at the nontrivial central vertex.
The five actions of $C^2,D^2,E^2,F^2,(CE)^2$ are exactly the
symplectic transvections of coefficient two, including the
transvection in direction $c+e$. Exhaustive multiplication of the
resulting four-by-four matrices modulo three gives $51840$ distinct
matrices. It also checks that
\[
                    (CE)^2 F^2 D^2 D^2 C^2
\]
acts on this subspace by a matrix of exact order five.
The matrices, group order, and word are retained in the JSON
certificate and exact output alongside the program. Any finite
$\{2,3\}$ total quotient dominating $F_6/N$ would make this action
factor through a finite $\{2,3\}$ group, contradicting that order.
The normalizer argument is exactly the one in
\statementref{A.CF6}. This finite computation verifies an instance
of the higher-dimensional mechanism in \statementref{A.CF8}; the
general proof of that statement does not rely on the computation.
\end{proof}

\subsection{Direct characteristic-positive proof, including proper wild covers}

\begin{claim}{A.CF12}{Proved}{Relative cuspidal exactness for arbitrary prime sets containing a non-residue prime}
Let \(C\) be a smooth proper geometrically connected curve of genus at
least two over an algebraically closed field of characteristic \(p>0\).
Let \(\Sigma\) be any set of primes and choose
\(r\in\Sigma\setminus\{p\}\). In the definitions of A.CF2, now formed
with maximal pro-\(\Sigma\) groups even when \(p\in\Sigma\), the
relative cuspidally pro-\(r\) configuration sequence is exact. The
boundary comparison and arithmetic transport of A.CF4 remain valid.
More generally, the finite-detection proof works for a constant proper
compactification with disjoint labelled sections in characteristic
\(p\), without a simultaneous lifting hypothesis.
\end{claim}
\begin{proof}
We replace the specialization step by a direct theorem whose arithmetic
base group is retained in full. In the opened text of
\citep[Proposition~2.7]{Stix2005}, the data are a connected locally noetherian scheme \(S\),
a smooth proper \(n\)-pointed genus-\(g\) curve over \(S\),
\(2-2g-n<0\), and a set \(L\) of primes invertible on \(S\).
Writing \(C^0\) for the complement of its sections, the asserted
sequence is
\[
 1\longrightarrow \pi_1(C^0_{\bar s})^{(L)}
  \longrightarrow \pi_1^{(L)}(C^0)
  \longrightarrow \pi_1(S)\longrightarrow1.             \tag{A.CF.12}
\]
The middle term is the quotient which pro-\(L\)-completes the fibre
kernel and retains the full group \(\pi_1(S)\). The proposition also
gives its continuous conjugation map to the automorphism group of the
fibre. Its proof constructs that map using the diagonal section after
base change by \(C^0\); the centrelessness of the hyperbolic pro-\(L\)
fibre proves injectivity. Thus the quoted statement is exactly a
relative completion assertion, and does not assert exactness after
pro-\(\Sigma\) completion of the total group.

Take any finite proper connected Galois \(\Sigma\)-cover
\(C'\to C\) with group \(A\), allowing \(p\mid\#A\).
The finite étale map is still a cover of smooth proper curves. Put
\[
 V=(C'\times C')\setminus\bigcup_{a\in A}\Gamma_a,
 \qquad S=C',\qquad L=\{r\}.
\]
Projection to the second factor has constant proper compactification
and disjoint labelled sections. Its geometric fibre has finitely
generated nonabelian free pro-\(r\) group \(P\).
The exact residue sequence for its mod-\(r\) homology has boundary
submodule generated by the labelled punctures, with their one sum
relation, and quotient \(H_1(C',\F_r)\). Both modules have trivial
base monodromy: the labels are fixed, the ground field is algebraically
closed, and the compactification is constant. Consequently the image
on \(P/\Phi(P)\) is contained in the additive group of shears from
the quotient module to the boundary module. It is a finite
\(r\)-group.

Compose \(\pi_1(V)\) with the conjugation map supplied by (A.CF.12).
The kernel of
\(\operatorname{Aut}(P)\to\operatorname{Aut}(P/\Phi(P))\)
is pro-\(r\), as proved in A.CF1. The image of this composite is
therefore pro-\(r\), so the homomorphism factors through
\(\pi_1(V)^{(r)}\). On the fibre it is the inner-automorphism
embedding of \(P\), which is faithful because \(P\) is centreless.
This proves the required finite-detection injection directly in
characteristic \(p\).

The remainder of A.CF2 only takes finite quotients and normal cores.
The resulting groups are an \(r\)-group by \(A\times A\), hence
are \(\Sigma\)-groups also when \(A\) has a \(p\)-part. All the
relative quotient maps and kernel identifications are unchanged.
For the boundary, proper covers remain étale along the diagonal,
whereas the new inertia has \(r\)-power order and is tame since
\(r\ne p\). The same diagonal root-gerbe construction and local
intersection argument from A.CF4 therefore apply. The arithmetic
assertion follows by retaining the entire absolute Galois quotient,
exactly as there.

For the asserted constant-compactification variant, first take a
common finite proper base cover which trivializes the pullbacks of
the chosen finite proper fibre cover along the finitely many section
maps. After this cover the lifted punctures are disjoint and labelled.
The same trivial-module residue sequence, Stix proposition, and
normal-core argument apply. No lifting of the family to characteristic
zero enters this proof.
\end{proof}
